\TieudeT{PHONG TỎA VẬT LÍ}
\subsubsection*{PHẦN I. Thí sinh trả lời câu hỏi từ câu 1 đến câu 18. Mỗi câu hỏi thí sinh chỉ chọn một phương án}
\setcounter{ex}{0}%mac
%\loai{An toàn khi sử dụng thiết bị điện}
\begin{ex}%book
Quy tắc về an toàn nào sau đây là \textbf{sai} khi sử dụng điện?
\choice
{\True Đảm bảo che chắn những cơ quan trọng yếu của cơ thể}
{Lắp đặt vị trí cầu dao, cầu chì, công tắc, ổ điện đúng nơi quy định}
{Tránh xa nơi điện thế nguy hiểm}
{Không dùng tay ướt hoặc nhiều mồ hôi khi sử dụng dây điện}	
\haidong{6}\end{ex}

\begin{ex}%book
Trong bài thực hành có sử dụng mạch điện nhưng khi lắp ráp xong mạch điện, báo cáo giáo viên phụ trách rồi cắm vào nguồn điện nhưng mạch không vào điện thì học sinh cần
\choice
{\True ngắt mạch điện ra khỏi nguồn sau đó kiểm tra mạch điện và nguồn điện}
{kiểm tra lại mạch điện}
{ngắt mạch điện ra khỏi nguồn}
{kiểm tra nguồn điện}	
\haidong{6}\end{ex}

%\loai{An toàn khi sử dụng thiết bị quang}
\begin{ex}%book
Phát biểu nào sau đây là \textbf{sai}? Khi sử dụng các thiết bị quang học cần chú ý
\choice
{sử dụng các thiết bị nhẹ nhàng}
{lau chùi cẩn thận thiết bị trước khi sử dụng}
{bảo quản nơi khô thoáng, tránh ẩm mốc}
{\True nhúng qua nước sôi để khử trùng}
\haidong{6}\end{ex}
%\loai{An toàn khi làm việc với chất phóng xạ}



%\newpage
%\loai{An toàn khi sử dụng thiết bị nhiệt và thủy tinh}

%\loai{Nhận biết biển báo}
\begin{ex}%book
\immini[thm]{Biển báo bên mang ý nghĩa gì?
\choice
{Bình chữa cháy}
{Chất độc môi trường}
{\True Bình khí nén áp suất cao}
{Dụng cụ dễ vỡ}}{\begin{tikzpicture}[draw=Mapcolor, very thick,>=stealth',transform shape,scale=0.12]
%\node at (0,0){\includegraphics[scale=0.4]{Hinh10}};
\fill[yellow,rounded corners=1mm] (0,7)--++(-60:14)--++(180:14)--cycle;
\draw[rounded corners=1mm,line width=3pt] (0,7)--++(-60:14)--++(180:14)--cycle;
\fill[rotate=37,rounded corners=1.5pt]  (-6.22,-0.15) rectangle (1.54,-1.22);
\fill[rotate=37,rounded corners=0.5pt]  (-5.92,-0.43) rectangle (2.46,-0.84);
\end{tikzpicture}}	
\haidong{6}\end{ex}

\begin{ex}%book
\immini[thm]{Biển báo bên là biển báo gì?
\choice
{\True Chất độc sức khỏe}
{Chất ăn mòn}
{Chất độc môi trường}
{Chất dễ cháy}}{
\begin{tikzpicture}[draw=Mapcolor, very thick,>=stealth',transform shape,scale=0.12]
%\node at (0,0.5){\includegraphics[scale=0.6]{Hinh17}};
\fill[yellow,rounded corners=2mm] (0,7)--++(-60:14)--++(180:14)--cycle;
\draw[rounded corners=2mm,line width=3pt] (0,7)--++(-60:14)--++(180:14)--cycle;
\fill[]  plot[smooth, tension=.4] coordinates {(-3.08,-1.7) (-3.04,-1.44) (-2.74,-1.38) (-2.18,-1.7) (2.06,-3.38) (3.06,-3.6) (3.14,-3.92) (2.86,-4.2) (2.72,-4.68) (2.18,-4.32) (1.62,-3.92) (-2.08,-2.46) (-2.32,-2.4) (-2.86,-2.5) (-3,-2.04) (-3.08,-1.7)};

\fill[] (-3.12,-3.96) .. controls (-3.18,-3.64) and (-2.94,-3.48) .. (-2.4,-3.44) .. controls (-0.84,-2.96) and (0.8,-2.28) .. (2.8,-1.2) .. controls (3,-1.14) and (3.1,-1.38) .. (3.02,-1.68) .. controls (2.92,-1.82) and (3,-1.92) .. (3.16,-2.06) .. controls (3.22,-2.34) and (2.94,-2.46) .. (2.46,-2.32) .. controls (2,-2.28) and (1.78,-2.42) .. (-0.26,-3.28) .. controls (-1.38,-3.74) and (-2.28,-4.1) .. (-2.2,-4.34) .. controls (-2.4,-4.64) and (-2.72,-4.58) .. (-2.68,-4.38) .. controls (-2.68,-4.12) and (-2.86,-4) .. (-3.12,-3.96);

\begin{scope}[xscale=-1]
\fill[] (-3.12,-3.96) .. controls (-3.18,-3.64) and (-2.94,-3.48) .. (-2.4,-3.44) .. controls (-0.84,-2.96) and (0.8,-2.28) .. (2.8,-1.2) .. controls (3,-1.14) and (3.1,-1.38) .. (3.02,-1.68) .. controls (2.92,-1.82) and (3,-1.92) .. (3.16,-2.06) .. controls (3.22,-2.34) and (2.94,-2.46) .. (2.46,-2.32) .. controls (2,-2.28) and (1.78,-2.42) .. (-0.26,-3.28) .. controls (-1.38,-3.74) and (-2.28,-4.1) .. (-2.2,-4.34) .. controls (-2.4,-4.64) and (-2.72,-4.58) .. (-2.68,-4.38) .. controls (-2.68,-4.12) and (-2.86,-4) .. (-3.12,-3.96);
\end{scope}
\fill[] (-2.22,1.08) .. controls (-1.46,3.52) and (1.34,3.6) .. (2.22,1.46) .. controls (2.32,0.9) and (2.16,0.18) .. (1.94,-0.38) .. controls (1.84,-0.54) and (1.84,-0.66) .. (1.94,-0.72) .. controls (2.22,-0.9) and (1.8,-1.32) .. (1.46,-1.3) .. controls (1.26,-1.1) and (0.92,-1.14) .. (0.98,-1.26) .. controls (1.02,-1.58) and (1.04,-1.78) .. (0.98,-1.96) .. controls (0.38,-2.18) and (-0.28,-2.22) .. (-0.94,-1.98) .. controls (-0.96,-1.74) and (-0.98,-1.54) .. (-0.98,-1.26) .. controls (-1.2,-1.16) and (-1.54,-1.42) .. (-1.76,-1.2) .. controls (-2.18,-0.94) and (-2.14,-0.74) .. (-1.9,-0.72) .. controls (-1.88,-0.52) and (-1.84,-0.5) .. (-1.92,-0.38) .. controls (-2.24,0.2) and (-2.34,0.7) .. (-2.22,1.08);

\fill[white] (-1.64,0.16) .. controls (-1.66,0.78) and (-0.64,0.96) .. (-0.32,0.42) .. controls (0.02,-0.4) and (-1.64,-0.9) .. (-1.64,0.16);
\fill[white] (-0.08,-0.24) .. controls (-0.2,-0.48) and (-0.64,-1.08) .. (-0.18,-1.12) .. controls (-0.08,-0.94) and (-0.02,-0.62) .. (-0.08,-0.24);
\begin{scope}[xscale=-1]
\fill[white] (-1.64,0.16) .. controls (-1.66,0.78) and (-0.64,0.96) .. (-0.32,0.42) .. controls (0.02,-0.4) and (-1.64,-0.9) .. (-1.64,0.16);
\fill[white] (-0.08,-0.24) .. controls (-0.2,-0.48) and (-0.64,-1.08) .. (-0.18,-1.12) .. controls (-0.08,-0.94) and (-0.02,-0.62) .. (-0.08,-0.24);
\end{scope}

\end{tikzpicture}
}

\haidong{6}\end{ex}

\begin{ex}%book
\immini[thm]{Biển báo bên mang ý nghĩa gì?
\choice
{Lối thoát hiểm}
{Cảnh báo tia laser}
{\True Lưu ý cẩn thận}
{Cảnh báo vật sắc, nhọn}}{\begin{tikzpicture}[draw=Mapcolor, very thick,>=stealth',transform shape,scale=0.12]
\fill[yellow,rounded corners=2mm] (0,7)--++(-60:14)--++(180:14)--cycle;
\draw[rounded corners=2mm,line width=3pt] (0,7)--++(-60:14)--++(180:14)--cycle;
\node at (0,0){\fontsize{287}{16}\bf !};
\end{tikzpicture}}

\haidong{6}\end{ex}
\begin{ex}%book
Kí hiệu $AC$ hoặc dấu "$\sim$" mang ý nghĩa gì?
\choice
{Cực dương}
{Dòng điện một chiều}
{Cực âm}
{\True Dòng điện xoay chiều}
\haidong{6}
\end{ex}

\begin{ex}%book
Trên các thiết bị thí nghiệm và trong phòng thí nghiệm thì kí hiệu nào sau đây cảnh báo có tia laser?
\choice
{\True 	\includegraphics[width=3cm]{images/2_11}}
{% TODO: \usepackage{graphicx} required
\includegraphics[width=3cm]{images/2_12}
}
{% TODO: \usepackage{graphicx} required
\includegraphics[width=3cm]{images/2_13}
}
{\includegraphics[width=3cm]{images/2_4}}
\haidong{6}\end{ex}



\begin{ex}%book
Trong các biển báo ở hình dưới, ở trong phòng thí nghiệm. Biển cảnh báo nào cảnh báo chất phóng xạ?
\begin{center}
\includegraphics[width=0.9\linewidth]{images/2_3}
\end{center}	
\choice
{Hình1}
{Hình 3}
{\True Hình 2}
{Hình 4}

\haidong{6}\end{ex}

\begin{ex}%book
\immini[thm]{Biển báo bên mang ý nghĩa gì?
\choice
{Cảnh báo tia laser}
{\True Nơi nguy hiểm về điện}
{Lưu ý cẩn thận}
{Cẩn thận sét đánh}}{\begin{tikzpicture}[draw=Mapcolor, very thick,>=stealth',transform shape,scale=0.12]
%\node at (0,0.5){\includegraphics[scale=0.6]{Hinh18}};
\fill[yellow,rounded corners=2mm] (0,7)--++(-60:14)--++(180:14)--cycle;
\draw[rounded corners=2mm,line width=3pt] (0,7)--++(-60:14)--++(180:14)--cycle;
\fill (-0.42,3.38) -- (-1.46,-1.38) -- (0.28,-0.44) -- (0.2,-2.5) -- (-0.12,-2.24) -- (0.26,-3.88) -- (1.22,-2.42) -- (0.78,-2.58) -- (1.56,0.94) -- (-0.44,-0.26) -- (0.94,3.4)--cycle;
\end{tikzpicture}}	
\haidong{6}\end{ex}

\begin{ex}%book
\immini[thm]{Biển cảnh báo ở hình bên liên quan đến
\choice
{chất ăn mòn}
{chất độc môi trường}
{chất độc sức khoẻ}
{\True chất dễ cháy}}{	\includegraphics[width=2cm]{images/2_9}}

\haidong{6}\end{ex}
%\loai{Nguy cơ nguy hiểm cho người sử dụng}

%\loai{Nguy cơ hỏng thiết bị đo}
\begin{ex}%book
Giới hạn đo của Ampe kế trong hình bên dưới là bao nhiêu?
\immini{		\choice[1]
{\True 3 A}
{$0.6\ A$}
{$0.02\ A$}
{$0.1\ A$}}{\includegraphics[width=4cm]{images/2_18}}

\haidong{6}\end{ex}

\begin{ex}%book
Phát biểu nào sau đây là \textbf{sai}? Để tránh nguy cơ hỏng thiết bị đo điện cần
\choice
{chọn chức năng và thang đo phù hợp}
{cắm dây đo vào chốt cắm phù hợp với thang đo}
{\True tùy ý cắm dây đo vào chốt cắm mình thích}
{không để các thiết bị đo ở gần nước}

\haidong{6}\end{ex}
%\loai{Nguy cơ cháy nổ và xử lí}

\begin{ex}%book
Khi phòng thực hành có đám cháy nơi có các thiết bị điện đang hoạt động thì thao tác nào sau đây là \textbf{sai}?
\choice
{Đưa toàn bộ các hoá chất, chất dễ cháy ra khu vực an toàn}
{Tổ chức thoát nạn}
{Ngắt toàn bộ hệ thống điện}
{\True Sử dụng nước dập đám cháy}	
\haidong{6}\end{ex}
%\loai{Các quy tắc an toàn}

\begin{ex}%book
Hoạt động nào sau đây \textbf{không} được làm sau khi kết thúc giờ thí nghiệm?
\choice
{Vệ sinh sạch sẽ phòng thí nghiệm}
{Sắp xếp gọn gàng các thiết bị và dụng cụ thí nghiệm}
{Bỏ chất thải thí nghiệm vào nơi quy định}
{\True Để các thiết bị nối với nguồn điện giúp duy trì năng lượng}
\haidong{6}\end{ex}

\begin{ex}%book
Quy tắc nào sau đây \textbf{không} đảm bảo an toàn trong phòng thực hành?
\choice
{Đọc kĩ hướng dẫn sử dụng thiết bị}
{Tuân thủ sự hướng dẫn của giáo viên hướng dẫn}
{Tắt công tắc nguồn thiết bị trước khi cắm điện và sau khi tháo điện}
{\True Tiếp xúc với nơi có cảnh báo nguy hiểm về điện}	
\haidong{6}\end{ex}

\begin{ex}%book
Khi tiến hành thí nghiệm, cần phải
\choice
{\True tuân theo các quy tắc an toàn của phòng thí nghiệm, hướng dẫn của giáo viên}
{tự đề xuất các quy tắc thí nghiệm để có thể tiến hành thí nghiệm nhanh nhất}
{thảo luận nhóm để thống nhất quy tắc riêng của nhóm, có thể bỏ qua quy tắc an toàn của phòng thí nghiệm}
{tiến hành thí nghiệm với thời gian ngắn nhất, không cần tuân thủ các quy tắc của phòng thí nghiệm}
\haidong{6}\end{ex}

\begin{ex}%book
Phát biểu nào sau đây là \textbf{sai}? Cần tuân thủ các biển báo an toàn trong phòng thực hành nhằm mục đích
\choice
{tránh được các tổn thất về tài sản nếu làm theo hướng dẫn}
{\True tạo ra nhiều sản phẩm mang lại lợi nhuận}
{hạn chế các trường hợp nguy hiểm như: đứt tay, ngộ độc,}
{chống cháy, nổ}

\haidong{6}\end{ex}



\subsubsection*{PHẦN 2. Thí sinh trả lời từ câu 1 đến câu 4. Trong mỗi ý a), b), c), d) ở mỗi câu, thí sinh chọn đúng hoặc sai.}
\setcounter{ex}{0}
\begin{ex}%Tailieuchuan
Trong phòng thực hành Vật lí, người học cần tuân thủ các quy tắc an toàn để tránh nguy hiểm khi sử dụng thiết bị điện và các nguồn sáng.
\choiceTF[t]
{Khi cắm điện vào nguồn điện ta cần thực hiện như hình bên trái} 
{\True Khi sử dụng dây điện ta không để dây như hình bên vì không an toàn} 
{Khi rút điện vào nguồn điện ta cần thực hiện như hình bên phải} 
{Chiếu tia laser mắt nhìn trực tiếp sẽ gây nguy hiểm cho mắt} 
\loigiai{ 
\begin{itemchoice}
\itemch
\itemch
\itemch
\itemch
\end{itemchoice}
}
\end{ex}
\begin{ex}%Tailieuchuan
Khi làm việc với các nguồn phóng xạ, cần tuân thủ nghiêm ngặt các nguyên tắc an toàn bức xạ.
\choiceTF[t]
{Thời gian tiếp xúc với nguồn phóng xạ không ảnh hưởng đến sức khỏe} 
{\True Để đảm bảo an toàn cần tăng khoảng cách từ người đến nguồn phóng xạ} 
{\True Để đảm bảo an toàn cần đảm bảo che chắn các cơ quan trọng yếu của cơ thể} 
{\True Nguồn phóng xạ khi bị rò rỉ ra môi trường sẽ gây tác hại cho: sức khỏe con người, môi trường nước, môi trường không khí...} 
\loigiai{ 
\begin{itemchoice}
\itemch
\itemch
\itemch
\itemch
\end{itemchoice}
}
\end{ex}
\begin{ex}%Tailieuchuan
Phóng xạ không chỉ gây hại mà còn có nhiều ứng dụng hữu ích trong đời sống và sản xuất.
\choiceTF[t]
{\True Được ứng dụng để điều trị bệnh ung thư} 
{Được ứng dụng trong nông nghiệp là làm giảm sức đề kháng của các loài sâu bọ} 
{\True Được ứng dụng trong khảo cổ để xác định niên đại của 1 mẫu vật} 
{\True Được ứng dụng để tìm vị trí rò rỉ nước trong các đường ống} 
\loigiai{ 
\begin{itemchoice}
\itemch
\itemch
\itemch
\itemch
\end{itemchoice}
}
\end{ex}

\begin{ex}%Tailieuchuan
Khi sử dụng các thiết bị thí nghiệm trong phòng thực hành Vật lí, người học cần nắm rõ các nguyên tắc an toàn.
\choiceTF[t]
{\True Các thiết bị đun nóng có thể gây bỏng với người sử dụng} 
{Các thiết bị đun nóng bằng thủy tinh sẽ không bị vỡ nên rất an toàn với người sử dụng} 
{\True Đường truyền tia sáng bị sai lệch là do các thiết bị quang học bị mốc, xước} 
{Khi các thiết bị quang học bị nứt thì đường truyền tia sáng không bị thay đổi} 
\loigiai{ 
\begin{itemchoice}
\itemch
\itemch
\itemch
\itemch
\end{itemchoice}
}
\end{ex}















\TieudeT{PHONG TỎA VẬT LÍ 10}
\subsubsection*{PHẦN I. Thí sinh trả lời câu hỏi từ câu 1 đến câu 18. Mỗi câu hỏi thí sinh chỉ chọn một phương án}
\setcounter{ex}{0}

\begin{ex}%book
Để sử dụng an toàn thiết bị đo điện khi sử dụng cần
\choice
{không chọn đúng thang đo, nhầm lẫn thao tác}
{chọn đúng thang đo, nhầm lẫn thao tác}
{không chọn đúng thang đo, thực hiện đúng thao tác}
{\True chọn đúng thang đo, thực hiện đúng thao tác}	
\haidong{6}\end{ex}
\begin{ex}%book
Phát biểu nào sau đây là \textbf{sai} khi nói về những quy tắc an toàn khi làm việc với phóng xạ?
\choice
{\True Mang áo phòng hộ và không cần đeo mặt nạ}
{Tăng khoảng cách từ ta đến nguồn phóng xạ}
{Đảm bảo che chắn những cơ quan trọng yếu của cơ thể}
{Giảm thời gian tiếp xúc với nguồn phóng xạ}

\haidong{6}%<MyLT1>
\end{ex}
\begin{ex}%book
Khi sử dụng các thiết bị nhiệt và thủy tinh trong phòng thí nghiệm Vật lí chúng ta cần lưu ý điều gì?
\choice
{Tiến hành thí nghiệm không cần quan sát vì tin tưởng vào dụng cụ phòng thí nghiệm}
{Quan sát các kí hiệu trên thiết bị, đặc điểm của các dụng cụ thí nghiệm, có thể dùng dụng cụ này thay thế cho dụng cụ khác}
{\True Quan sát kĩ các kí hiệu trên thiết bị, đặc điểm của dụng cụ thí nghiệm, chức năng của dụng cụ}
{Có thể sử dụng mọi ống thủy tinh trong phòng thí nghiệm vào tất cả các thí nghiệm}
\haidong{6}\end{ex}
\begin{ex}%book
Kí hiệu DC hoặc dấu "$-$" mang ý nghīa gì?
\choice
{Cực dương}
{Dòng điện xoay chiều}
{\True Dòng điện một chiều}
{Cực âm}
\haidong{6}\end{ex}	
\begin{ex}%book
Kí hiệu nào trên dụng cụ thí nghiệm dùng điện, thực hành sau đây bị \textbf{sai}?
\choice
{Input là đầu vào}
{Màu đỏ là cực dương}
{Output là đầu ra}
{\True AC là dòng điện một chiều}

\haidong{6}\end{ex}
\begin{ex}%book
Kí hiệu "Input (I)" mang ý nghĩa gì?
\choice
{Cực dương}
{\True Đầu vào}
{Đầu ra}
{Cực âm}

\haidong{6}\end{ex}
\begin{ex}%book
Kí hiệu \includegraphics[width=0.07\linewidth]{images/2_14}
mang ý nghĩa gì?
\choice
{Dụng cụ dễ vỡ}
{Không được phép bỏ vào thùng rác}
{\True Tránh ánh nắng chiếu trực tiếp}
{Dụng cụ đặt đứng}

\haidong{6}\end{ex}	
\begin{ex}%book
\immini[thm]{Kí hiệu ở hình bên được in trên các thiết bị thí nghiệm mô tả điều gì?
\choice
{Lưu ý cẩn thận}
{\True Từ trường}
{Dụng cụ đặt đứng}
{Dụng cụ dễ vỡ}}{	\includegraphics[width=2cm]{images/2_4}}

\haidong{6}\end{ex}
\begin{ex}%book
Biển báo 
\includegraphics[width=0.05\linewidth]{images/2_6}
mang ý nghĩa gì?
\choice
{\True Cần mang đồ bảo hộ cơ thể, chống hóa chất, chống nước}
{Cần mang kính bảo vệ mắt khỏi những hóa chất độc hại và đảm bảo thị lực của người trong phòng thí nghiệm}
{Bao tay chống hóa chất, chống khuẩn}
{Cần đeo mặt nạ phòng độc}

\haidong{6}\end{ex}
\begin{ex}%book
\immini[thm]{Kí hiệu biển báo bên có ý nghĩa gì?
\choice
{Nhiệt độ cao}
{Cảnh báo tia laser}
{\True Cảnh báo phóng xạ}
{Nguy hiểm về điện}}{\includegraphics[width=2cm]{images/2_2}}

\haidong{6}\end{ex}	
\begin{ex}%book
Khi làm thí nghiệm, thực hành Vật lí, nếu phát hiện người bị điện giật, ta phải làm gì đầu tiên?
\choice
{Gọi cấp cứu}
{Gọi người đến sơ cứu}
{\True Ngắt nguồn điện}
{Đưa người bị điện giật ra khỏi khu vực có điện}	
\haidong{6}\end{ex}
\begin{ex}%book
\immini[thm]{Một thiết bị trong phòng thí nghiệm có các đặc tính như hình. Thiết bị này sẽ cho điện áp đầu ra là
\choice
{dòng điện xoay chiều với điện áp $240\ V$}
{\True dòng điện xoay chiều với điện áp $12\ V$}
{dòng điện một chiều với điện áp $240\ V$}
{dòng điện một chiều với điện áp $12\ V$}}{	\includegraphics[width=4cm]{images/2_19}}

\haidong{6}%<MyLT1>	
\end{ex}
\begin{ex}%book
Khi có sự cố chập cháy dây điện trong khi làm thí nghiệm ở phòng thực hành, điều ta cần làm trước tiên là
\choice
{\True ngắt nguồn điện}
{dùng nước để dập tắt đám cháy}
{dùng $CO_2$ để dập đám cháy nếu chẳng may lửa cháy vào quần áo}
{thoát ra ngoài}
\haidong{6}\end{ex}
\begin{ex}%book
Khi gặp sự cố mất an toàn trong phòng thực hành, học sinh cần
\choice
{\True báo cáo ngay với giáo viên trong phòng thực hành}
{tự xử lí và không báo với giáo viên}
{nhờ bạn xử lí sự cố}
{tiếp tục làm thí nghiệm}

\haidong{6}\end{ex}	

\begin{ex}%book
Để đảm bảo an toàn trong phòng thực hành cần thực hiện nguyên tắc nào dưới đây?
\choice
{Có thể nhận biết hóa chất bằng cách ngửi hóa chất}
{Mang đồ ăn vào phòng thực hành}
{Làm thí nghiệm theo sự hướng dẫn của bạn bè trong lớp}
{\True Đọc kĩ nội quy và thực hiện theo nội quy phòng thực hành}	
\haidong{6}\end{ex}	
\begin{ex}%book
Hoạt động nào sau đây \textbf{không} thực hiện đúng quy tắc an toàn trong phòng thực hành?
\choice
{Đeo găng tay khi làm thí nghiệm}
{Không ăn uống, đùa nghịch trong phòng thí nghiệm}
{\True Để hóa chất không đúng nơi quy định sau khi làm xong thí nghiệm}
{Làm thí nghiệm theo sự hướng dẫn của giáo viên}	
\haidong{6}\end{ex}	
\begin{ex}%book
Hành động nào \textbf{không} tuân thủ quy tắc an toàn trong phòng thực hành?
\choice
{Trước khi cắm, tháo thiết bị điện, ta tắt công tắc nguồn trước}
{Trước khi làm thí nghiệm với bình thủy tinh, cần kiểm tra bình có bị nứt vỡ hay không}
{Bố trí dây điện gọn gàng}
{\True Dùng tay không để làm thí nghiệm}	
\haidong{6}\end{ex}	
\begin{ex}%book
Phát biểu nào sau đây là \textbf{sai} khi nói về nguy cơ gây mất an toàn khi sử dụng thiết bị thí nghiệm Vật lí trong phòng thực hành?
\choice
{Nhờ giáo viên kiểm tra mạch điện trước khi bật công tắc nguồn điện}
{\True Thực hiện những thí nghiệm với nhiệt độ cao mà không dùng đồ bảo hộ}
{Bố trí dây điện gọn gàng, kiểm tra mạch có điện bằng bút thử điện}
{Trước khi cắm hoặc rút phích cắm của thiết bị điện, phải tắt công tắc nguồn}
\haidong{6}\end{ex}






\subsubsection*{PHẦN 2. Thí sinh trả lời từ câu 1 đến câu 4. Trong mỗi ý a), b), c), d) ở mỗi câu, thí sinh chọn đúng hoặc sai.}
\setcounter{ex}{0}
\begin{ex}
Hình ảnh là ảnh chụp nhãn thông số trên thiết bị chuyển đổi điện áp được sử dụng cho máy tính xách tay.
\begin{center}
\includegraphics[width=8 cm]{images/2025_04_24_fe017539a8fd00ee4a76g-1(2)}
\end{center}
\choiceTF[t]
{Chức năng của thiết bị là thay đổi cường độ dòng điện}
{\True Thiết bị được sử dụng cho dòng điện xoay chiều}
{Khi thiết bị hoạt động bình thường, đầu ra là dòng điện xoay chiều có hiệu điện thế $12\ V$}
{\True Khi thiết bị không còn sử dụng được, cần phân loại riêng để xử lí}
\loigiai{
\begin{itemchoice}
\itemch 
\itemch 
\itemch 
\itemch 
\end{itemchoice}
}
\end{ex}
\begin{ex}
Hình ảnh là ảnh chụp đồng hồ đo điện đa năng hiện số.
\immini{\choiceTF[t]
{Đồng hồ chỉ dùng để đo hiệu điện thế}
{Đồng hồ đang ở chế độ đo hiệu điện thế một chiều}
{\True Khi đo cường độ dòng điện một chiều, giới hạn đo của thiết bị là $10\ A$}
{\True Khi sử dụng chế độ đo dòng một chiều, que đo nối với cực âm sẽ được cắm vào cổng COM}}{
\includegraphics[width=2.7 cm]{images/2025_04_24_fe017539a8fd00ee4a76g-1}}
\loigiai{
\begin{itemchoice}
\itemch 
\itemch 
\itemch 
\itemch 
\end{itemchoice}
}
\end{ex}

\begin{ex}
Ở một phòng thí nghiệm có treo một số biển báo sau.
\begin{center}
\includegraphics[width=12 cm]{images/103136}
\end{center}
\choiceTF[t]
{\True Biển báo $a$ là biển báo cấm}
{Biển báo $b$ và $c$ thuộc cùng một loại biển báo}
{\True Cần đeo găng tay khi thực hiện thí nghiệm}
{Khi xảy ra hỏa hoạn, không được tự ý dập lửa}
\loigiai{
\begin{itemchoice}
\itemch 
\itemch 
\itemch 
\itemch 
\end{itemchoice}
}
\end{ex}

\begin{ex}
Trong bài thực hành đo cường độ dòng điện của bóng đèn, bạn A phải sử dụng máy biến áp, bóng đèn, dây dẫn và đồng hồ đo điện đa năng. Để quá trình thực hiện thí nghiệm diễn ra được an toàn thì bạn A cần phải thực hiện thao tác.
\choiceTF[t]
{\True Kiểm tra cẩn thận tất cả các thiết bị trước khi sử dụng}
{Bật công tắc nguồn thiết bị ngay sau khi hoàn thành việc lắp ráp mạch điện}
{Thực hiện thí nghiệm gần khu vực có chứa cồn}
{\True Ngay lập tức sử dụng bình cứu hỏa ở nơi gần nhất để dập tắt đám cháy khi có hỏa hoạn xảy ra}
\loigiai{
\begin{itemchoice}
\itemch 
\itemch 
\itemch 
\itemch 
\end{itemchoice}
}
\end{ex}


